\subsection{交叉乘积}

设 K 是一个体，并设 G 是一个群，我们用 e 表示它的单位元素。设 L 是一个交换 K-代数，并设 \(\lambda\) 是从 G 到 K-代数 L 的自同构群的一个同态。对 G 中的任何 g，用 \(\tau(g)\) 表示 L 的乘法群 \(L^{*}\) 的、由 \(\lambda(g)\) 诱导的自同构。

设 \(\mathcal{E} = (\Gamma, \iota, \pi)\) 是一个 \(\tau\)-扩张（群 \(G\) 被 \(L^*\) 的）。用下式定义群 \(L^*\) 在集合 \(L \times \Gamma\) 上的一个右作用

\[
(\beta , \gamma). \alpha = (\beta \alpha , \iota (\alpha) ^ {- 1} \gamma)\tag{20}
\]

对 \(\alpha \in \mathrm{L}^*\)、\(\beta \in \mathrm{L}\) 与 \(\gamma \in \Gamma\)。我们用 E 表示 \(\mathrm{L}^*\) 在 \(\mathrm{L} \times \Gamma\) 中的轨道的集合，并用 \([\beta; \gamma]\) 表示偶 \((\beta, \gamma)\) 的轨道。于是按构造我们有关系

\[
[ \beta \alpha ; \gamma ] = [ \beta ; \iota (\alpha) \gamma ]\tag{21}
\]

对 \(\alpha\in\mathrm{L}^{*}\)、\(\beta\in\mathrm{L}\) 与 \(\gamma\in\Gamma\)。

给定 \(\beta\)（L 中的）与 \(\gamma\)（\(\Gamma\) 中的），我们用 \(\gamma\beta\) 表示 \(\beta\) 在 L 的自同构 \(\lambda\circ\pi(\gamma)\) 下的变换；我们有关系

\[
{ } ^ { \gamma } ( { } ^ { \gamma } { } ^ { \prime } \beta ) = { } ^ { \gamma } { } ^ { \gamma } { } ^ { \prime } \beta   , \quad { } ^ { \gamma } ( \beta + \beta ^ { \prime } ) = { } ^ { \gamma } \beta + { } ^ { \gamma } \beta ^ { \prime }   , \quad \text {and} \quad { } ^ { \gamma } ( \beta \beta ^ { \prime } ) = { } ^ { \gamma } \beta   { } ^ { \gamma } \beta ^ { \prime }\tag{22}
\]

对 \(\gamma,\gamma'\)（\(\Gamma\) 中的）与 \(\beta,\beta'\)（L 中的）。E 上存在一个由关系式刻画的合成律

\[
[ \beta ; \gamma ] [ \beta^ {\prime}; \gamma^ {\prime} ] = [ \beta^ {\gamma} \beta^ {\prime}; \gamma \gamma^ {\prime} ].\tag{23}
\]

事实上，只需验证：若我们把 \(\beta\)、\(\gamma\)、\(\beta'\) 与 \(\gamma'\) 分别换成 \(\beta\alpha\)、\(\iota(\alpha)^{-1}\gamma\)、\(\beta'\alpha'\) 与 \(\iota(\alpha')^{-1}\gamma'\)（对所有 \(\alpha\)、\(\alpha'\)（\(L^*\) 中的）），右端不变。这由对 \(\mathcal{E}\) 应用 VIII, 第 285 页的公式 (1) 立即得出，该公式也可以写成 \(\gamma\iota(\alpha) = \iota(\gamma\alpha)\gamma\)（对 \(\gamma \in \Gamma\) 与 \(\alpha \in L^*\)）。由 (22) 中的公式，装备了该律的集合 E 是一个以 {[}1; e{]} 为单位元素的幺半群。

由于 \(\pi \circ \iota\) 是取常值 \(e\) 的映射，存在一个从 E 到 G 的映射 \(\widetilde{\pi}\)，使得我们有

\[
\widetilde {\pi} ([ \beta ; \gamma ]) = \pi (\gamma)\tag{24}
\]

对 \(\beta \in \mathrm{L}\) 与 \(\gamma \in \Gamma\)。

设 \(g\) 是 \(\mathrm{G}\) 的一个元素，并设 \(\mathrm{E}_g = \widetilde{\pi}^{-1}(g)\)。若 \(\gamma_0\) 是 \(\pi^{-1}(g)\) 的一个固定元素，则映射 \(\beta \mapsto [\beta; \gamma_0]\) 是从 \(\mathrm{L}\) 到 \(\mathrm{E}_g\) 的一个双射，借助它我们将把 K-向量空间结构转移到 \(\mathrm{E}_g\) 上（该结构是 \(\mathrm{L}\) 上的）。由于 \(\pi^{-1}(g)\)

由形如 \(\iota(\alpha)\gamma_{0}\) 的元素组成（其中 \(\alpha\) 取遍 \(L^{*}\)），且我们有

\[
[ \beta ; \iota (\alpha) \gamma_ {0} ] = [ \beta \alpha ; \gamma_ {0} ],
\]

所以 \(E_{g}\) 上的向量空间结构不依赖于 \(\gamma_{0}\) 的选取。

设 \(g\) 与 \(g'\) 是 \(\mathrm{G}\) 的元素；由公式 (22) 与 (23)，\(\mathrm{E}\) 上的合成律经限制诱导一个从 \(\mathrm{E}_g \times \mathrm{E}_{g'}\) 到 \(\mathrm{E}_{gg'}\) 的 K-双线性映射。于是向量空间 \(\mathrm{P} = \bigoplus_{g \in \mathrm{G}} \mathrm{E}_g\) 装备了一个结合且含幺代数的结构，其乘法诱导上述从 \(\mathrm{E}_g \times \mathrm{E}_{g'}\) 到 \(\mathrm{E}_{gg'}\) 的双线性映射（对所有 \(g\) 与 \(g'\)，均在 \(\mathrm{G}\) 中）。代数 \(\mathrm{P}\) 称为 \(\mathrm{L}\) 与 \(\mathcal{E}\) 的交叉乘积，记作 \(\mathbf{A}[\mathcal{E};\mathrm{L}]\)；它的单位元素是元素 \([1;e]\)（\(\mathrm{E}_e\) 的）。

令

\[
u (\beta) = [ \beta ; e ]\tag{25}
\]

对 L 中的 \(\beta\)。于是 \(u: \mathrm{L} \to \mathbf{A}[\mathcal{E};\mathrm{L}]\) 是 K-代数的一个单同态。由 (23)，对每个 \(\gamma \in \Gamma\)，元素 \([1;\gamma]\) 在 \(\mathbf{A}[\mathcal{E},\mathrm{L}]\) 中可逆，且映射 \(v: \Gamma \to \mathbf{A}[\mathcal{E},\mathrm{L}]^*\)（它把 \(\gamma\) 映到 \([1;\gamma]\)）是一个单群同态。同态 \(u\) 与 \(v\) 称为典范的。我们有关系

\begin{enumerate}
\def\labelenumi{(\arabic{enumi})}
\setcounter{enumi}{25}
\tightlist
\item
\end{enumerate}

\[
u (\alpha) = v (\iota (\alpha)),\tag{27}
\]

\[
u (\gamma \beta) = v (\gamma) u (\beta) v (\gamma) ^ {- 1},\tag{28}
\]

\[
[ \beta ; \gamma ] = u (\beta) v (\gamma)
\]

对 \(\alpha \in \mathrm{L}^*\)、\(\beta \in \mathrm{L}\) 与 \(\gamma \in \Gamma\)。

反过来，代数 \(A[E;L]\) 具有下述泛性质。

命题 12. —— 设 B 是一个 K-代数，\(u' : L \to B\) 是一个 K-代数同态，而 \(v' : \Gamma \to B^{*}\) 是一个群同态。我们假设下列关系

\begin{enumerate}
\def\labelenumi{(\arabic{enumi})}
\setcounter{enumi}{28}
\tightlist
\item
\end{enumerate}

\[
u ^ {\prime} (\alpha) = v ^ {\prime} (\iota (\alpha)),\tag{30}
\]

\[
u ^ {\prime} (\gamma \beta) = v ^ {\prime} (\gamma) u ^ {\prime} (\beta) v ^ {\prime} (\gamma) ^ {- 1}
\]

对 \(\alpha \in \mathrm{L}^*\)、\(\beta \in \mathrm{L}\) 与 \(\gamma \in \Gamma\) 成立。则存在唯一的代数同态 \(f\)（从 \(\mathbf{A}[\mathcal{E};\mathrm{L}]\) 到 \(\mathrm{B}\)），使得我们有 \(u' = f \circ u\) 与 \(v' = f \circ v\)。

为证明 \(f\) 的唯一性，注意向量空间 \(\mathbf{A}[\mathcal{E};\mathbf{L}]\)（体 \(\mathrm{K}\) 上的）由形如 \([\beta ;\gamma ] = u(\beta)v(\gamma)\) 的元素组成的集合生成。

若同态 \(f' : A[\mathcal{E}; L] \to B\) 满足 \(f' \circ u = u'\) 与 \(f' \circ v = v'\)，则它把 \([\beta; \gamma]\) 映到 \(u'(\beta)v'(\gamma)\)，因此与 f 重合。

由公式 (29)，我们有

\[
u ^ {\prime} (\beta \alpha) v ^ {\prime} (\iota (\alpha) ^ {- 1} \gamma) = u ^ {\prime} (\beta) v ^ {\prime} (\gamma)\tag{31}
\]

对 \(\alpha\)（\(\mathrm{L}^*\) 中的）、\(\beta\)（\(\mathrm{L}\) 中的）与 \(\gamma\)（\(\Gamma\) 中的）。于是由 \(\mathrm{E}\) 的定义，存在一个映射 \(f_0: \mathrm{E} \to \mathrm{B}\)，使得 \(f_0([\beta; \gamma]) = u'(\beta)v'(\gamma)\)。由公式 (23) 与 (30)，我们有 \(f_0(xx') = f_0(x)f_0(x')\)，其中 \(x\) 与 \(x'\) 在 \(\mathrm{E}\) 中。\(f_0\) 在 \(\mathrm{E}_g\) 上的限制对每个元素 \(g\)（\(\mathrm{G}\) 的）都是 K-线性的；因此，存在唯一的 K-线性映射 \(f\)（从 \(\mathbf{A}[\mathcal{E}; \mathrm{L}]\) 到 \(\mathrm{B}\)），它与 \(f_0\) 在 \(\mathrm{E}\) 上重合。映射 \(f\) 是一个满足 \(u' = f \circ u\) 与 \(v' = f \circ v\) 的代数态射。

注. —— 设 \(\sigma: G \to \Gamma\) 是映射 \(\pi\) 的一个截面。对 G 中的每个 g 令 \(\varepsilon_{g} = v(\sigma(g))\)，并用 \(c_{\sigma}\) 表示与 \(\sigma\) 相伴的 2-闭上链（VIII, 第 296 页）。特别地，我们有

\[
\varepsilon_ {g} \varepsilon_ {g ^ {\prime}} = u (c _ {\sigma} (g, g ^ {\prime})) \varepsilon_ {g g ^ {\prime}}\tag{32}
\]

对所有 \(g, g' \in G\)。此外，让我们借助同态 u 把 L 与 \(A[E; L]\) 的一个子代数等同起来。于是 \(A[E; L]\) 的每个元素都可以唯一地写成 \(\sum_{g \in G} a_g \varepsilon_g\)，其中 \((a_g)\) 是 L 的元素的一个有限支撑族。\(A[E; L]\) 中的乘法可以用公式

\[
\left(\sum_ {g} a _ {g} \varepsilon_ {g}\right) \left(\sum_ {g} b _ {g} \varepsilon_ {g}\right) = \sum_ {g} d _ {g} \varepsilon_ {g}\tag{33}
\]

其中

\[
d _ {g} = \sum_ {h h ^ {\prime} = g} a _ {h} (\lambda (h) \cdot b _ {h ^ {\prime}}) c _ {\sigma} (h, h ^ {\prime}).\tag{34}
\]

若扩张 \(\mathcal{E}\) 是半平凡的，则我们可以选取一个截面 \(\sigma\)（\(\pi\) 的），它是从 G 到 \(\Gamma\) 的一个群态射；于是闭上链 \(c_{\sigma}\) 是取常值 1 的，且公式 (34) 化简为

\[
d _ {g} = \sum_ {h h ^ {\prime} = g} a _ {h} \lambda (h) \cdot b _ {h ^ {\prime}}.\tag{35}
\]

设 \(K'\) 是体 K 的一个扩张。用 \(L'\) 表示 \(K'\)-代数 \(L_{(K')}\)，并对 G 中的任何 g，用 \(\lambda'(g)\) 表示自同构 \(\lambda(g)_{(K')}\)（由 \(\lambda(g)\) 在 \(L_{(K')}\) 上诱导的）。此外，我们用 \(\tau'(g)\) 表示 \(L'^*\) 的、由 \(\lambda'(g)\) 诱导的自同构。最后，设 h 是从 \(L^*\) 到 \(L'^*\) 的把 x 映到 \(x \otimes 1\) 的同态。设 \(\mathcal{E}' = (\Gamma', \iota', \pi')\) 是 E 的扩张被 h 的直像（VIII, 第 289 页）。设 \(A[\mathcal{E}', L']\) 是交叉乘积 \(K'\)-代数（\(\mathcal{E}'\) 与 \(L'\) 的），并设 \(u': L' \to A[\mathcal{E}', L']\) 与 \(v': \Gamma' \to A[\mathcal{E}', L']^*\) 是典范同态。

命题 13. —— 存在唯一的 \(\mathrm{K}'\)-代数同构 \(\varphi\)（从 \(\mathbf{A}[\mathcal{E};\mathrm{L}]_{(\mathrm{K}')}\) 到 \(\mathbf{A}[\mathcal{E}';\mathrm{L}']\)），它满足关系

\[
u ^ {\prime} (h (\beta)) = \varphi (1 \otimes u (\beta))\tag{36}
\]

对 \(\beta \in \mathrm{L}\)，以及

\[
v ^ {\prime} (h (\gamma)) = \varphi (1 \otimes v (\gamma))\tag{37}
\]

对 \(\gamma\in\Gamma\)。

证明由诸构造立即得出。

同构 \(\varphi\) 称为典范的。
% semantic-fragments: legacy-input-seam
\relax{}
